%%%%%%%%%%%%%%%%%%%%%%%%%
\section{\label{sec:checklist}Best Practices Checklist for Deriving MC Trial Acceptance}
%%%%%%%%%%%%%%%%%%%%%%%%%

The checklist to obtain MC trial acceptance probabilities in a way that minimizes the possibility of error is provided on the next page.
The acceptance probabilities are derived in a step-by-step manner that is closely related to the simulation algorithms.
By starting from the list of instructions to fully describe a forward and reverse trial, a researcher is able to obtain straightforwardly the required transition probabilities, and identify cases where irreversible changes must be rejected, in a systematic fashion.
We apply this checklist in subsequent sections to derive a variety of unbiased and biased MC trials in various ensembles.
This process of generating the needed acceptance probabilities was inspired by and built upon a set of MC lecture notes \cite{kofke_lecture9_2024} and MC ideal gas simulation tests of various theoretically-based ensemble predictions \cite{hatch_theory_2024}.

\begin{Checklists*}[p!]
\begin{checklist}{Best Practices Checklist for Deriving MC Trial Acceptance}
%\begin{checklist}{\label{sec:checklist}Best Practices Checklist for Deriving MC Trial Acceptance}
\begin{itemize}
\item
  \textbf{Write the MC trial algorithm as an ordered list of all steps required to fully describe both the forward and reverse trials.}

\item
  \textbf{Check detailed balance.}
  \begin{itemize}
    \item Ensure all forward trial steps transition from the old state to a new state.
    \item Consider many different cases, including edge cases such as the largest or smallest values.
    \item Ensure all steps conducted in the reverse trial are capable of returning from the new state to the exact same old state where the forward trial began.
  \end{itemize}

\item
  \textbf{Determine if each step is deterministic, stochastic or needs to be broken into multiple steps.}
  \begin{itemize}
    \item Stochastic steps may choose a particle, position, volume or model parameter.
    \item Stochastic steps make one choice from multiple options.
    \item Stochastic steps often use a random number generator (seeded for reproducibility).
    \item Deterministic steps make no choices.
    \item If more than one choice is made in one step, break that step into multiple steps.
  \end{itemize}

\item
  \textbf{Determine the probability for each step.}
  \begin{itemize}
    \item Determine a probability for choosing a discrete variable (E.g. if the step requires choosing a particle from among $N$ particles, then the probability is $1/N$).
    \item Determine a probability for choosing a continuous variable (E.g., if the step requires choosing a position in a volume $V$, then the probability is $\diff\mathbf{r}/V$).
    \item If the step is deterministic, then the probability is 1.
    \item Determine the forward trial probability, $\alpha_{o \rightarrow n}$, from the old to the new state.
    \item Determine the reverse trial probability, $\alpha_{n \rightarrow o}$, from the new to the old state.
    \item Obtain the total probability of the forward trial as the product of the forward probabilities for each step.
    \item Repeat with the previous step, for the total reverse probability.
  \end{itemize}

\item
  \textbf{Determine the ensemble of the trial.}
  \begin{itemize}
    \item List all possible variables that could change during each step (Note if the particle positions, number of particles, volume or model parameters change).
    \item From the list of all variables that could change, determine the ensemble.
    \item Use $NVT$ ensemble microstate probabilities if only particle positions change.
    \item Use $NPT$ ensemble microstate probabilities if only particle positions and volume change.
    \item Use (semi-)$\mu VT$ ensemble microstate probabilities if the number of particles changes.
    \item Use Gibbs ensemble microstate probabilities if multiple systems are coupled.
    \item Use Expanded ensemble microstate probabilities if temperature or model parameter change (beyond the scope of this article).
  \end{itemize}

\item
  \textbf{Use Eq.~\ref{eq:detailed_balance} to obtain the trial acceptance probability.}
  \begin{itemize}
    \item The ensemble of the trial determines the ratio of microstate probabilities.
    \item The probabilities of each step of the trial implementation determines the ratio of the forward and reverse transition probabilities.
  \end{itemize}

\item
  \textbf{Test the implementation of the trial acceptance probability.}
  \begin{itemize}
    \item Verify that the acceptance probability leads to simulation results that agree with theory and trusted sources.
    \item Start with the simplest possible simulations that are fast, such as ideal gases with tractable theoretical expressions.
  \end{itemize}

\item
  \textbf{Document your derivation or software implementation for later reference.}
  \begin{itemize}
    \item Include documentation and references for future readers (e.g., yourself) to understand the derivation.
  \end{itemize}

\end{itemize}
\end{checklist}
\end{Checklists*}
